%%%%%
%
\section{Humidit� sp�cifique � partir de l'humidit� relative}

\itemstyle{Op\'erateur~:} {\sf Humspec\_Hu}
\index{Humspec\_Hu}

%%%%%OBJET%%%%
\itemstyle{Objet~:}
  Calcul d'une humidit� sp�cifique (rapport entre la
masse de vapeur d'eau et la masse d'air \textbf{humide} :
$q = \frac{m_v}{m_v+m_a})$
� partir de l'humidit� relative et des param�tres d'�tat
(pression, temp�rature).

%%%%%ARGUMENTS%%%%
\bigskip
%%%%%ENTREES%%%%
\begin{tabarg}
\multicolumn{3}{l}{\titretab {Entr\'ees}}\\
\hline
$Hu$ & Vecteur, {\tt float} & Humidit� relative (\%)\\
$T_s$ & Vecteur, {\tt float} & Temp�rature (Kelvin ou Celsius)\\
$P_s$ & Vecteur, {\tt float} & Pression (hPa)\\
\end{tabarg}

\medskip
%%%%%SORTIES%%%%
\begin{tabarg}
\multicolumn{3}{l}{\titretab {Sorties}}\\
\hline
$q$ & Vecteur, {\tt float} & Humidit� sp�cifique (g/kg)\\
\end{tabarg}


%%%%%FORMULE%%%%
\itemstyle{Formule}

%%Humspec_Hu
\[
  e = \frac{Hu}{100} . e_w(T_s)
\]
puis
\[
  q = 1000 . \frac{R_a}{R_v} . \frac{e}{P_s - e (1-\frac{R_a}{R_v})}
\]



%%%%%AUTEUR%%%%
\itemstyle{Auteur~:}
CNRM/GMEI/TRAMM

%%%%%REFERENCE%%%%
\itemstyle{R\'ef\'erences~:}
D�finitions des humidit�s (relatives, absolues, sp�cifiques)
selon la table 4.3 des tables m�t�orologiques internationales
de l'OMM \cite{OMM-188.TP.94}
